\documentclass[11pt,a4paper]{article}
\usepackage[T1]{fontenc}
\usepackage{lmodern,amsmath,amssymb,amsthm,booktabs,microtype,hyperref}
\pdfmapfile{+lm.map}
\usepackage[margin=27mm]{geometry}
\setlength{\emergencystretch}{2em}
\newcommand{\Q}{\mathbb Q}\newcommand{\Z}{\mathbb Z}
\newcommand{\Tr}{\operatorname{Tr}}\newcommand{\rad}{\operatorname{rad}}
\newcommand{\N}{\operatorname{N}}
\newtheorem{theorem}{Theorem}\newtheorem{lemma}[theorem]{Lemma}
\newtheorem{proposition}[theorem]{Proposition}
\title{Audit and strengthening of the single-pole weak-$abc$ construction}
\author{Z diophantine in $\Q$ project}
\date{19 November 2026}
\begin{document}\maketitle
\begin{abstract}
The construction of Gabe Conant with Astra~6.2 combines the two contact-height branches
of the earlier thirteen-variable construction in one pole predicate. We
verify the rational scalar, the ordered local preparation, and the finite
algebra guard logic. The resulting eleven-variable theorem is valid under
its expressly retained arithmetic premises. Full auxiliary elliptic
charts, ten further automatically nonzero factors, and a birational
normalization of the pole square class improve the two displayed degree
bounds to $32542081$ and $19366273$. The sufficient fixed input shift
remains ineffective. A split-contact formulation explains the full range
$1<\kappa<3/2$ and permits odd positive scale exponents. None of these
arguments establishes a ten- or nine-variable defining polynomial.
\end{abstract}

\paragraph{Source credit.}
The supplied constructions and notes audited here are by Gabe Conant with
Astra~6.2, in the research bundles
\emph{H10-quartic-constructions} and
\emph{H10-eleven-variable-constructions}.
The audit, corrections, and further improvements stated in this paper are
subsequent project work. The original source files and their provenance
remain separately preserved.

\section{Scope and precise statements}
Let $W_\kappa$ assert that there is a constant $K_\kappa$ with
$c\le K_\kappa\rad(abc)^\kappa$ for every coprime positive integer triple
$a+b=c$. We retain the arithmetic inputs of the project's v22 and v28
constructions: the stated elliptic group and index-chart identification,
the one-point trace-height estimate, the simultaneous prime-pattern
selection with its specified finite congruences and order compatibility,
and the native pole obstruction and prepared converse. The latter includes
the density and finite-character avoidance used in v28; mere existence of
one rational point on one pole fiber is not substituted for that input.
These hypotheses are unchanged by the present transformations.

In particular put
\[
 K=\Q(\sqrt2),\quad D=11-6\sqrt2,\quad E:y^2=x^3-Dx,
 \quad 2E(K)=\Z P_0,
\]
where
\[
 P_0=\left((57-28\sqrt2)/8,(-660+417\sqrt2)/32\right).
\]
The retained root chart represents $A=[n]P_0$ with $n\ne0$ and
\(T=\Tr x(A)\), and supplies constants $c_0,j_0$ such that
\(n^2\le c_0 h(T)+j_0\).

\begin{theorem}\label{thm:main}
Assume these retained inputs and $W_\kappa$.
For either row below there is an explicitly specified family of integer
polynomials in eleven variables, including the input, whose rational input
projection equals $\Z$ for every sufficiently large fixed odd integer $B$.
\[
\begin{array}{c|c|c|r|r}
\kappa&e&k&\text{Conant--Astra bound}&\text{improved degree bound}\\\hline
121/100&6&6&33344385&32542081\\
16/15&4&5&19861377&19366273
\end{array}
\]
The support bounds for the improved polynomials are respectively
$<2^{235}$ and $<2^{228}$. The rank of the finite algebra is $128$.
The coefficient files at $B=3$ are reproducible fixtures; no claim that
$B=3$ exceeds the sufficient defining threshold is made.
\end{theorem}
The families are specified by the native rational DAG, integral finite
algebra, and strengthened determinant receipt. All counts include the
input. The degree statements are upper bounds, not claims about an
expanded determinant having exactly that degree.

\section{Eight split contacts and the scalar inequality}
The exact identity
\[
 (w+4)(w+20)^3(w-5)^2-(w-4)(w-20)^3(w+5)^2=108w^5
\]
is verified by expansion. The degree-six forms obtained by homogenizing
have no common projective zero. Their three fibers consist of
$0,\infty,\pm4,\pm5,\pm20$; the fixed transformation $t=1/(w-21)$ gives
\[
 \Theta=\{0,-1/21,-1/17,-1/25,-1/16,-1/26,-1,-1/41\}.
\]
Write $s=8$ and $G(t)=\prod_{\theta\in\Theta}(t-\theta)$.
The cited LMFDB orbit is a primary reference for the Belyi map, while the
identity and coprimality needed here are independently exact.

Enlarge a fixed set $S$ by $2,3,7,17$, the fixed arithmetic exceptions,
and primes of contact denominators, nonzero contact differences, and
$1-\theta$. For a primitive numerator-denominator pair representing $t$,
let $N(t)$ be the product of the eight primitive homogeneous contact forms.
The bounded resultant gcd in the degree-six identity and $W_\kappa$ give
\[
 \log\rad N(t)\ge(6/\kappa)h(t)-O(1),\qquad
 \log|N(t)|\le8h(t)+O(1).
\]
Hence the product of primes outside $S$ dividing $N(t)$ exactly once
has logarithm at least $(12/\kappa-8)h(t)-O(1)$. Indeed each exponent
$r\ge2$ has $2-r\le0$, so this simple-prime product is at least
$\rad(N)^2/|N|$ up to a fixed factor. At such a prime the finite contact
product $G(t)$ has valuation one. No prime dividing the denominator of
$t$ divides a good finite contact form.

Set
\[
 \begin{gathered}
 \beta=X^2+B,\qquad a=\beta^e+T^e,\qquad x=a-z,\quad y=a+z,\\
 P=G(x),\qquad Q=G(y),\qquad J=(x-1)(y-1),
 \end{gathered}
\]
\[
 F=P^2+Q^3,\quad
 \mathcal D=(P^3+Q^5)(P+Q^k)(P^k+Q),\qquad
 u=\frac{J^{s(k-1)}PQF^2}{\mathcal D}.
\]
Initially protect $P,Q,x-1,y-1$ and the original point charts. The four
contact sums are automatically nonzero by Lemma~\ref{lem:positive} below.
Let $r(u)$ be the product of primes outside $S$ with positive odd valuation
on $u$.

\begin{lemma}\label{lem:scalar}
For $k\ge2$ and $s=8$,
\[
 \log r(u)\ge\left(12/\kappa-8-8/k\right)
                  \max(h(x),h(y))-O(1).
\]
The constants do not depend on $B,n,z$.
\end{lemma}
\begin{proof}
Choose the slot of greater height and its simple primes. Discard those at
which the other integral slot's contact product has valuation at least
$k$. The lost logarithm is at most $(8/k)$ times the other height plus a
fixed constant, by its primitive octic numerator bound.
At a remaining prime with both slots integral let $m=v(P),n'=v(Q)$.
One equals one and the other is less than $k$. If the other equals zero,
$F$ and all three factors of $\mathcal D$ are units, and
$v(u)=1+s(k-1)v(J)$, positive odd. If both are positive, $J$ is a unit and
all relevant minima are strict, so
\[
 v(u)=2\min(2m,3n')-\min(3m,5n')=1.
\]
For $m=1$ the two minima are $2,3$; for $n'=1,m\ge2$ they are $3,5$.
If instead $v(P)=1$ and $t_0=v(y)<0$, then $v(Q)=s t_0$ and
$v(J)=t_0$. Direct substitution gives
\[
 s(k-1)t_0+1+s t_0+6s t_0-(k+6)s t_0=1.
\]
The reversed mixed case gives
$s(k-1)t_0+1+s t_0+4s t_0-(k+4)s t_0=1$.
This includes denominator primes of the moving scale; none was discarded
merely because the scale has a pole.
\end{proof}

Let $E_h=h(1:\beta:T)$. The source binary-power argument gives
\[
 h(a)\ge(e/\kappa-3)E_h-O(1).
\]
Here $\beta>0$ and $T=2g_0^2+4g_1^2>0$ on the root chart, so the proof
also applies to odd positive $e$: its use of evenness was solely to ensure
positive summands and $H^e\le A_0^e+B_0^e\le2H^e$.
Since $x+y=2a$, $\max(h(x),h(y))\ge h(a)/2-O(1)$. Consequently
\begin{equation}\label{eq:lambda}
 \log r(u)\ge\lambda E_h-O(1),\qquad
 \lambda=\tfrac12(12/\kappa-8-8/k)(e/\kappa-3).
\end{equation}
For the two rows of Theorem~\ref{thm:main}, $\lambda$ is respectively
$8374/14641$ and $99/160$, strictly greater than $1/2$.
In particular $r(u)>1$ for all such scalar specializations when $B$ is
sufficiently large, because $E_h\ge\log B$. This conclusion is obtained
before any balancing denominator or pole equation is used.

\section{Ten more automatic factors}
\subsection{Four point guards and the character-extraction order}
There are three rational-base point norm factors and three ordinate-norm
factors in the unchanged native charts. Write $\mathsf W_A=\N(k_0)$
for the root ordinate norm, where $k_0^2=g^4-D$, and write
$\mathsf W_H,\mathsf W_U$ for the full auxiliary ordinate norms.
Retain $\mathsf W_H,\mathsf W_U$, but remove the norms of
$g,x(H),x(U)$ and the norm factor for $\mathsf W_A$.

The recovery occurs partly before and partly after obtaining a rational
character. For every rational base $g\in K$, the equation $g^4=D$
is impossible: taking norms would give $(\N g)^4=49$, whereas no rational
number has fourth power $49$. Consequently
\[
 \mathsf W_A^2=\N(g^4-D)\ne0
\]
holds at every rational base, independently of any rational character.
For either full auxiliary abscissa $h\in K$, $\N h=0$ would imply $h=0$;
then $h^3-Dh=0$ and $\mathsf W_H=0$, contrary to its retained norm
factor. The auxiliary abscissa norms are therefore also recovered before
character extraction.

All three point blocks are now etale. In the v71 primitive-point proof,
the only necessary rational-base hypothesis is $a^2-2b^2\ne0$ for its
ordinate radicand $a+b\sqrt2$. It implies both $2a^2-b^2\ne0$ and
$-4a^3+6ab^2+3b^3\ne0$; the latter follows from the Eisenstein polynomial
$4T^3-6T-3$ when $b\ne0$. Thus the six point-prefix quadratic generators
remain nonzero and separating, and the terminal root may still vanish
under the v40 argument. Apply the determinant and Galois separation to
obtain a rational character. At that character $k_0\in K$; if $g=0$,
then $k_0^2=-D$, impossible because $D=11-6\sqrt2$ is totally positive.
This recovers the last removed point norm, $\N g$.

No rational expression with a removed denominator has been evaluated
before this recovery. The cleared integral multiplication algebra and
its separator polynomials exist on those base loci. The argument first
recovers the exact hypotheses used for character extraction, and only
then evaluates the native point and detector expressions.

\subsection{Four positive contact sums and two odd-valuation factors}
\begin{lemma}\label{lem:positive}
Let $G$ be monic of even degree with every root in $[-1,0]$, and suppose
$x+y=2a>4$. Then
\[
 P^2+Q^3>0,\quad P^3+Q^5>0,\quad P+Q^k>0,\quad P^k+Q>0
 \qquad(k\ge2).
\]
\end{lemma}
\begin{proof}
Outside $[-1,0]$ the even-degree product $G$ is nonnegative. If $G(t)<0$,
then $t\in[-1,0]$ and every $|t-\theta|\le1$, so $|G(t)|\le1$.
At least one of $x,y$ is at least $a>2$, and its contact product exceeds
one. If both contact products are nonnegative every displayed expression
is positive. If one is negative, its absolute value is at most one and
the positive opposite product exceeds one; the same conclusion follows
in each expression, including both parities of $k$.
\end{proof}

The following recovery order is essential. First recover the point units
as above and retain the individual contact and $J$ factors. Lemma~\ref{lem:positive}
recovers the four contact sums and makes $u$ well defined and nonzero.
Equation~\eqref{eq:lambda} then supplies a prime $p\notin S$ of positive
odd valuation. Therefore $u\ne-1$. It also recovers every balance
denominator as shown in the next section. The positive detector factors
then recover the norm and trace of their product. Finally the native pole
conic denominator is recovered from the same odd prime. No deleted
factor is used to justify its own nonvanishing.

\section{Full auxiliary points and one detector}
The root $A$ keeps $x=g^2,y=gk_0,k_0^2=g^4-D$. The auxiliary points use
\[
 H:(x,y)=(h,j),\ j^2=h^3-Dh,\qquad
 U:(x,y)=(v,l_0),\ l_0^2=v^3-Dv.
\]
Each element of $K$ here has two retained rational abscissa coordinates
and a rank-four finite ordinate algebra. On the guards these are precisely
the affine points of $E(K)$ with nonzero abscissa and ordinate.
For example if $h=h_0+h_1\sqrt2$, the radicand coefficients are
\[
 h^3-Dh=(h_0^3+6h_0h_1^2-11h_0+12h_1)
 +(3h_0^2h_1+2h_1^3+6h_0-11h_1)\sqrt2.
\]
The v71 rank-four multiplication table and two quadratic point stages
apply; the auxiliary stage degrees are $6,9$, instead of $8,12$.

Keep the six physical tests, with $f=1/x_E$, $\zeta=-x_E/y_E$,
$r=\zeta(U)/\zeta(H)$, and conjugation denoted by a bar:
\[
 \begin{gathered}
 C_1=f(H),\qquad C_2=f(U),\\
 C_3=\overline{f(H-P_0)},\qquad C_4=\overline{f(U-A)},\\
 C_5=r-\overline{\zeta(U-A)/\zeta(H-P_0)},\qquad C_6=r-X.
 \end{gathered}
\]
Set $c_i=C_i$ for $i\le4$ and $c_i=C_i^2$ for $i=5,6$, and put
\[
 t_i=\frac{c_i}{1+u c_i},\qquad
 W=\prod_{i=1}^6\left(1+(t_i/u)^2\right),\quad V=\Tr(W^{-1}),
 \quad b=\frac{4u}{(1+u)^2V^2}.
\]
For the first four tests the homogeneous implementations $R/(L+uR)$
retain the old nonzero $L$; no extraneous affine chart is introduced.

At every prime outside $\{2,7\}$, a nonzero abscissa of any $K$-rational
point of $E$ has even valuation. If its valuation is positive or negative,
the curve equation respectively has valuations $v(x)$ or $3v(x)$,
which must be even; valuation zero is already even. Thus every nonzero
$c_i$ has even valuation outside the fixed set. This argument replaces the
stronger global-square assertion, which is not available for unrestricted
auxiliary points.

At the positive odd prime obtained above, $1+u c_i$ cannot vanish: its
two terms have different valuations. It is consequently a nonzero element
of $K$ and invertible in both embeddings. At any real rational point
character, all factors of $W$ are at least one, so $\N(W)>0$ and
$V>0$. Hence their scalar norms and trace guards may be omitted after
this ordered recovery.

\section{The pole normalization and soundness}
Write the exact native square extraction as $b=d_0 t_0^2$. Let
$\delta=P+Q^k>0$, and replace it by
\[
 d=d_0/\delta^2,\qquad t=t_0\delta,\qquad b=dt^2.
\]
Use free rational parameters $\tau,s_0$, and a terminal quadratic root $q$:
\[
 e_0=\frac{\tau^2-17d}{\tau^2+51d},\qquad
 q^2=17d\bigl(s_0^2-17e_0(17-b)\bigr)
           \bigl(s_0^2+51e_0(17-b)\bigr).
\]
The old and new pole charts are isomorphic on the guarded source:
$\tau_{\rm old}=\delta\tau$ and $q_{\rm old}=\delta q$, with $s_0$ unchanged.
Thus both rational existence and density for avoiding finitely many
proper algebraic conditions transfer.

At the chosen odd prime, $v_p(d)\equiv v_p(b)\equiv v_p(u)\pmod2$.
The constants $3,17$ are units there. Hence $\tau^2+51d\ne0$, including
$\tau=0$, because otherwise an even valuation would equal an odd one.
This completes the six scalar and pole guard deletions in addition to
the four point guards. The terminal ordinate may
vanish; the retained pole obstruction explicitly covers that case.

For soundness take any prime outside $S$ with $\nu=v_p(u)>0$ odd.
The saturation identity at either place above $p$ is
\[
 v(t_i)=\max(v(c_i),-\nu),
\]
where a zero test is treated separately as a factor one. Parity makes
$v(t_i/u)\ne0$, so each detector factor has nonpositive valuation.
Put $q_*=\max_{\mathfrak p\mid p}v_{\mathfrak p}(W)\le0$.
Since $v_p(V)\ge-q_*$ and $b/u$ is a nonzero rational square, the native
negative-odd obstruction gives
\[
 0<v_p(b)\le\nu+2q_*.
\]
At a place attaining $q_*$ every detector factor has valuation
$> -\nu/2$, forcing $v(t_i)>3\nu/4$ and hence $v(C_i)>0$.

For arbitrary auxiliary points write $2H=[h']P_0$, $2U=[j']P_0$ and
$h=h'/2,j=j'/2$. These are locally integral at our odd primes.
The guarded logarithmic quotients are $j/h$ and $(j-n)/(h-1)$.
The formal comparison and the six positive tests imply
$Xh\equiv j$ and $X(h-1)\equiv j-n$ modulo $p$, with $X$ integral there.
Subtracting gives $X\equiv n$. This is exactly the v71 extension of the
original local implication; no integer label for $H$ itself is assumed.

If $X\ne n$, therefore $r(u)$ divides the nonzero reduced numerator of
$X-n$. As $2h(X)\le h(\beta)\le E_h$,
\[
 (\lambda-1/2)E_h\le\log(1+|n|)+O(1).
\]
The trace estimate $n^2\le c_0E_h+j_0$ bounds $E_h$ uniformly, contradicting
$E_h\ge\log B$ for every sufficiently large fixed $B$. This proves
soundness, including all sign possibilities for $u$.

\section{Completeness and its order of choices}
Fix a nonzero integer $n$, its actual root $A=[n]P_0$, and hence $a$.
Choose the inherited real and finite-place neighborhoods on the
multiplication family $U=[n]H$, avoiding the original chart poles.
Use the old $2E(K)$ subfamily for these prepared witnesses. The full
auxiliary charts preserve all geometric characters: if the old chart is
$x=g^2,y=gk_0$, set $h=g^2,j=gk_0$; multiplication by the guarded unit $g$
identifies the two square-root algebras in each embedding.

Before selecting $z$, enlarge the temporary finite set by the denominators
of $a$ and the contacts, their nonzero differences, the nonzero specialized
values at $(x,y)=(1,2a-1),(2a-1,1)$, and the three resultants of $F$ with
the factors of $\mathcal D$. The resultants are nonzero: a common complex
zero either has $P=Q=0$, or the absolute-value equations force
$|P|=|Q|=1$. In the latter case $|x|,|y|\le2$, since $|G(t)|>1$ for
$|t|>2$; in the former the same bounds hold. Both contradict $x+y=2a>4$.

At every fixed place, $u/z^8\to1$ as $|z|\to\infty$.
On the chosen point neighborhoods, $W\to1,V\to2$, and $bz^8\to1$.
This prepares local square classes and arbitrarily deep positive even
valuations at $2,3,17$, and $0<b<17$ over the reals.
Apply the retained prime theorem to the sixteen distinct rational
positions $\pm(a-\theta)$, with all finite constraints already fixed,
to choose $z=N/(d_1^2L^2)$ and distinct suitable prime numerators.
Refine the auxiliary index by the compatible CRT conditions while
preserving the earlier neighborhoods.

At a fresh prime exactly one contact has valuation one; the other contact
product and $J$ are units. Thus $v_p(u)=1$. In the successful field slot
the first four reciprocal abscissas have valuation at least two and the
last two squares also have valuation at least two. Consequently
$W\equiv1$ there. In the other slot $W^{-1}$ is congruent to zero or one
by the strict parity argument. Hence $V$ is a unit and $v_p(b)=1$.

At every unprepared, unselected integral slope prime, $P,Q$ are units.
A prime of $\mathcal D$ is not a prime of $F$ or $J$ by the prepared
specializations and resultants, so it gives $v_p(u)<0$; otherwise the
possible positive contributions from $J$ and $F$ are even. At a new pole
of $z$, $v_p(u)=8v_p(z)$ is even. Negative odd $\nu=v_p(u)$ gives
$v(t_i/u)\ge-2\nu>0$ in both slots, so $W\equiv1,V\equiv2$ and
$v_p(b)=-\nu>0$. Even valuations of $u$ give even valuations of $b$.
Thus the reduced denominator of $b$ is a square at every prime.
The native prepared converse now supplies the pole parameters, with the
retained density furnishing finite-character avoidance. The new square
class normalization is birational by the nonzero scalar $\delta$.

This verifies all prime classes, not only the freshly selected primes.
No choice of contact primes precedes a subsequently changing value of $a$.
The usual factor $X$ includes the integer zero independently.

\section{The integral determinant and exact checks}
The free coordinates are $X$, two root abscissa coordinates, four
auxiliary abscissa coordinates, $z$, and two pole parameters: ten in
all. The character encoder adds one. The three point blocks and the
terminal square block have ranks $4,4,4,2$, for total rank $128$.
For the seven quadratic generators use the positive polynomial envelopes
$M_i$ in the parent specification and the triangular separator
\[
 \phi_7=q_7/(4M_7),\qquad
 \phi_i=q_i(1+\phi_{i+1})/(4M_i),\qquad
 \Omega=\left(\prod_i4M_i\right)\phi_1.
\]
The first six generators are units on the retained point guards; the
last may vanish. The v40 proof separates distinct characters even on
this terminal zero stratum and makes $\Omega$ nonzero. Multiplication
in the integral free algebra is defined by the supplied coefficient
DAG and its cleared relations. If $G_0$ is the product of the retained
subalgebra norms, the polynomial is
\[
                 F_B=X\det\bigl(I-YG_0\operatorname{Mult}(\Omega)\bigr).
\]
A determinant zero makes $G_0\ne0$ and recovers a rational character by
Galois separation. The automatic factors are then recovered in the
order already proved. Conversely a prepared character gives rational
$Y=(G_0\Omega)^{-1}$. The coefficient of $X Y^0$ is one, so the integer
content is one.

For the two preferred specifications the independently recomputed guard
and separator degrees are
\[
 (g,w)=(490,253744),\qquad (418,150880).
\]
Thus $\deg F_B\le1+128(1+g+w)$ gives exactly the asserted upper bounds.
The support bound uses total pre-input degree $d=128(1+g+w)$ and
$Y$-degree at most $128$:
\[
 \#\operatorname{supp}(F_B)\le
 \binom{d+11}{11}-\binom{d-128+10}{11}.
\]
The fixture coefficient one-norm bounds are $2^{131679488}$ and
$2^{83631616}$ at $B=3$. Their scope is deliberately distinct from the
uniform-in-$B$ degree and support claims.

The independent checks reconcile the point and test expressions with
v71, reconstruct the entire scalar and detector circuit, verify
$b=dt^2$, recompute the subalgebra bounds, and verify the staged deleted
atom sets. Corrupted degree or guard receipts are rejected. These finite
checks certify the algebra and numerical ledger; they are not numerical
proofs of the external arithmetic inputs or an effective defining shift.

\section{Generalizations and remaining limits}
More generally, suppose a coprime homogeneous degree-$m$ triple
$A+B=C$ has completely split rational branch support of size $s$.
A fixed rational change of coordinate puts that support in $[-1,0]$.
The same argument gives simple-prime coefficient $\gamma=2m/\kappa-s$
and, for the shared scalar, the sufficient condition
\[
 \tfrac12(2m/\kappa-s-s/k)(e/\kappa-3)>1/2.
\]
For even $s$ the preceding proof applies unchanged. Thus any such split
triple with $1<\kappa<2m/s$ gives an eleven-variable family after choosing
$k$ and then $e$ sufficiently large. The degree-six triple gives
$1<\kappa<3/2$.
For odd $s$, the underlying guarded construction still works with odd
$k$ and locally square prepared $z$: then $s(k-1)$ is even, and every
new pole of the square-denominator slope has even valuation. The four
positive-contact-sum deletions are not asserted for odd $s$, because a
monic odd-degree contact product is negative on an unbounded interval.
The resulting guard version is a generalization of the construction,
not of the specific degree receipts above.

The two displayed $(e,k)$ choices remain the best among the nearby
parameters checked with the same compiler, including the closest odd
competitors. Odd exponents can improve other fixed weak exponents; for
example $\kappa=7/5$ permits $e=11,k=22$ with
$\lambda=272/539>1/2$. No higher-degree completely split triple is claimed
from a limited database scan, and no variable elimination follows merely
from an auxiliary curve being everywhere locally soluble.

\begin{thebibliography}{9}
\bibitem{Coauthor} Gabe Conant with Astra~6.2, supplied research package,
\emph{H10-eleven-variable-constructions}, research note dated 18 November
2026, \texttt{note/weak-eleven.tex} and \texttt{counts/poles/}.
\bibitem{Project} J. Brox and contributors,
\emph{Polynomial presentations of the integers over the rationals},
commit \texttt{d3a8c223716a0fc0dbf5af4eb7998d640c21048d}.
\url{https://github.com/JoseBrox/diophantine-z-q_improvements}.
Relevant complete proofs: direct-input v22, finite algebras v28/v30,
terminal roots v40, full auxiliary points v71, even abscissa valuations v75.
\bibitem{Belyi} M. Musty, S. Schiavone, J. Sijsling and J. Voight,
\emph{A database of Belyi maps}, Open Book Series 2 (2019), 375--392.
The exact orbit used here is
\url{https://www.lmfdb.org/Belyi/6T16/5.1/3.2.1/3.2.1/a/}.
\end{thebibliography}
\end{document}
